% SPDX-License-Identifier: Apache-2.0
% covers: I2c
\datasheet{I2c}{A general I2C master on AXI-Lite: one piece of a transaction per command, on two open drain lines.}

\dsfacts{\code{lib/parts/src/i2c.rs}}{\code{txhdl\_parts::i2c::I2c}}%
{\code{//docs:examples}}

\subsection*{Function}

\code{I2cInit} of the HDMI part replays a fixed list of writes and
cannot read. This master is the other kind: a program says what one
piece of a transaction is and the master does it, so a driver can
address a device, write to it, turn the bus around and read from it.

One command is at most a start, one byte either way, and a stop. A
driver that wants to read a register writes the address with
\code{start|write}, then asks for \code{start|read|nack}, whose start
is the repeated start. Between the commands of one transaction the
lines are parked with the clock low, since a data line that moves
while the clock is high is a start or a stop to every device on the
bus.

\subsection*{Register map}

Offsets from the block's base.

\begin{center}\footnotesize
\begin{tabular}{@{}llp{0.52\textwidth}@{}}
\toprule
Offset & Word & Access \\
\midrule
\code{0x00} & \code{ctrl} & Read and write. Bits 0 to 15 the divider,
bit 16 the interrupt enable. \\
\code{0x04} & \code{cmd} & Write only; writing it starts the command.
Bit 0 start, 1 stop, 2 write, 3 read, 4 the answer after a read, bits
8 to 15 the byte to write. \\
\code{0x08} & \code{data} & Read only, the byte a read took. \\
\code{0x0c} & \code{state} & Read, and write ones to clear bits 1, 2
and 3. Bit 0 busy, 1 done, 2 not acknowledged, 3 arbitration lost. \\
\bottomrule
\end{tabular}
\end{center}

A quarter of a bit takes \code{div + 1} cycles, so the bus runs at the
system clock over \code{4 * (div + 1)}. From 100 MHz, 249 is
100 kbit/s and 62 is 400 kbit/s.

Bit 4 of \code{cmd} is set for a NACK, which is what a master sends
after the last byte it wants; a master that acknowledges the last byte
tells the device to send another.

\dstables{I2c}

\subsection*{The two lines}

Both are open drain, as I2C requires: a device pulls a line low and
nobody drives one high. \code{scl\_low} and \code{sda\_low} high pull
their line low, and a line nobody pulls is high through the board's
pull-up, so a board wrapper drives a pad from each output and reads
the pad back into \code{scl\_in} and \code{sda\_in}.

The master watches both inputs, and two things follow that a master
driving the lines itself could not do.

A device that holds the clock low is stretching it, and the master
waits for the line to rise rather than counting cycles. And a data
line that reads low while this master released it is another master
winning the bus, which sets arbitration lost and ends the command.

\subsection*{Behaviour}

\begin{itemize}
\item Writing \code{cmd} clears the three status bits, records what
was asked, and starts the run. \code{step} walks the transaction:
idle, the start, the eight bits, the acknowledge, the stop, the end.
\item \code{quarter} says which quarter of the bit the lines are in,
and \code{tick} counts the cycles of that quarter against the divider.
\item On a write, the acknowledge step reads the data line: low is the
device acknowledging, high sets the not-acknowledged bit.
\item On a read, the byte arrives in \code{shift} and is put in the
\code{data} register, which the word at \code{0x08} reads. The master
then drives the answer bit 4 asked for.
\item Finishing a command sets the done bit, which raises the
interrupt while the enable is set, and clears busy.
\item \code{held} records that a start has been sent and no stop has,
which is what parks the lines between the commands of one
transaction.
\end{itemize}

\subsection*{Verification}

\begin{itemize}
\item \code{//lib/parts:parts\_test} drives the master against the
device model in the same file.
\item \code{//lib/examples:ex\_i2c} addresses a device, writes to it,
turns the bus around and reads back, and the build simulates the
netlist against that run under nvc and Verilator.
\end{itemize}

\subsection*{Used by}

Nothing in the tree yet. The HDMI part still uses \code{I2cInit} for
its fixed setup; issue 148 notes that \code{I2cInit} could become a
program for this master instead, and that a board's EEPROM, which
often holds the Ethernet address, is what wants a master that can
read.

\subsection*{Limits}

Ten-bit addressing is not here: an address is a byte like any other,
so a driver writes what it likes, but nothing helps it. There is no
general call handling, no clock synchronisation with another master
beyond losing arbitration, and no timeout on a stretched clock, so a
device that holds the clock low for ever holds the master with it.
